% Context from chapters-tex/sec-optim-smooth.tex; for mathematical reference, not compilation.
\section{Gradient Descent Algorithm}
\label{sec-grad-desc-basic}


                                                                                  
\subsection{Steepest Descent Direction}

The Taylor expansion~\eqref{eq-grad-dfn} approximates $f$ near $x$ by an affine function:
\eq{
	f(z) = T_x(z) + o(\norm{x-z})
	\qwhereq
	T_x(z) \eqdef f(x) + \dotp{\nabla f(x)}{z-x}, 
}
Figure~\ref{fig-expansion-taylor} illustrates the approximation. First-order methods use $f$ through its local model $T_x$.


\begin{figure}[tbp]
\centering
\BookPanel{.485}{1}{tikz/smooth-taylor-line}\hfill
\BookPanel{.485}{2}{tikz/smooth-taylor-plane}
\caption{Panels 1 and 2: first-order Taylor approximations in one and two dimensions.
	Panels 3 and 4: the gradient is orthogonal to the level set; the diagram illustrates the proof.}
\label{fig-expansion-taylor}
\end{figure}
\begin{figure}[tbp]
\BookContinuedFigure
\centering
\BookPanel{.485}{3}{tikz/smooth-level-set-tangent}\hfill
\BookPanel{.485}{4}{tikz/smooth-nested-level-sets}
\caption{Panels 1 and 2: first-order Taylor approximations in one and two dimensions.
	Panels 3 and 4: the gradient is orthogonal to the level set; the diagram illustrates the proof.}
\label{fig-expansion-taylor-continued}
\end{figure}

A nonzero gradient $\nabla f(x)$ points uphill locally, so $-\nabla f(x)$ is a descent direction. At a fixed point $x$, examine $f$ along the ray
\eq{ 
	\tau \in \RR^+ = [0,+\infty[ \longmapsto f(x-\tau \nabla f(x)) \in \RR.
}
If $f$ is differentiable at $x$, one has 
\eq{
	f(x-\tau \nabla f(x)) = f(x) - \tau \dotp{\nabla f(x)}{\nabla f(x)} + o(\tau)
		= f(x) - \tau \norm{\nabla f(x)}^2 + o(\tau).
}
If $\nabla f(x)=0$, then $x$ is stationary and is a global minimizer when $f$ is convex. Otherwise, for sufficiently small $\tau$,
\eq{
	f(x-\tau \nabla f(x)) < f(x)
} 
so the step from $x$ to $x-\tau \nabla f(x)$ decreases the objective.

\begin{rem}[Orthogonality to level sets]
	A level set of $f$ contains points with the same value of $f$. For $s \in \RR$, define
	\eq{
		\Ll_s \eqdef \enscond{x}{f(x)=s}.		
	}
	Assume $f$ is continuously differentiable near $x$ and $\nabla f(x)\neq0$. Then $x$ is a regular point of the level set $\Ll_{f(x)}$. The gradient is normal to its tangent space and points toward increasing values of $f$; see Figure~\ref{fig-expansion-taylor-continued}, panels 3 and 4.
	 
	Take a smooth curve through $x$ in $\Ll_s$, parameterized by $t \in \RR \mapsto c(t)$ with $c(0)=x$.
	  
	The function $h(t) \eqdef f(c(t))$ has constant value $h(t)=s$ because $c(t)$ lies in the level set. Thus $h'(t)=0$. Expanding at $t=0$ also gives
	\eq{
		h(t) = f(c(0) + t c'(0) + o(t)) = h(0)+t\dotp{c'(0)}{\nabla f(c(0))} + o(t)
	}	
	i.e. $h'(0) = \dotp{c'(0)}{\nabla f(x)}=0$, so that $\nabla f(x)$ is orthogonal to the tangent $c'(0)$ of the curve $c$, which lies in the tangent plane of $\Ll_s$  (as shown in Figure~\ref{fig-expansion-taylor-continued}, panel 3).  Since the curve $c$ is arbitrary, the whole tangent plane is thus orthogonal to $\nabla f(x)$.
\end{rem}

\begin{rem}[Local optimal descent direction]
	Assume $f$ is continuously differentiable near $x$ and $\nabla f(x)\neq0$. Among displacements of length $r$, every minimizing direction approaches the normalized negative gradient as $r\downarrow0$:
\eq{
	\frac{1}{r} \uargmin{\norm{u}=r} f(x + u) \overset{r \rightarrow 0}{\longrightarrow} -\frac{\nabla f(x)}{\norm{\nabla f(x)}}.
}
To verify the sign, let $u_r$ minimize $f(x+u)$ on the radius-$r$ sphere. Compare it with $-r\nabla f(x)/\norm{\nabla f(x)}$. The uniform first-order remainder gives
\eq{\dotp{\nabla f(x)}{u_r/r}\leq-\norm{\nabla f(x)}+o(1).}
Since $\norm{u_r/r}=1$, equality in the limiting Cauchy--Schwarz bound forces the stated direction.
\end{rem}


                                                                                  
\subsection{Gradient Descent}

Starting from $x_0 \in \RR^p$, gradient descent iterates
\eql{\label{eq-grad-desc}
	x_{k+1} \eqdef x_k - \tau_k \nabla f(x_k)
}
The parameter $\tau_k>0$ is the step size, or learning rate. A sufficiently small $\tau_k$ decreases $f$ at a nonstationary iterate. Choosing $\tau_k$ balances stable descent against progress per step.
 
One may fix $\tau_k=\tau$ or adapt $\tau_k$ at each iteration; see Figure~\ref{fig-gradesc}.


\begin{figure}[tbp]
\centering
\BookPanel{.485}{1}{tikz/smooth-gradient-step-size}\hfill
\BookPanel{.485}{2}{tikz/smooth-exact-line-search}
\caption{Effect of the step size $\tau$ on gradient descent (panel 1) and exact line search (panel 2).}
\label{fig-gradesc}
\end{figure}


\begin{rem}[Greedy choice]
Exact line search chooses the best step along the current direction, although computing it can be expensive:
\eq{
	\tau_k\in\uargmin{\tau\geq0} h(\tau) \eqdef f(x_k-\tau \nabla f(x_k)).
}
The line-search objective $h(\tau)$ is scalar. Its derivative follows from a Taylor expansion of $h$:
\eq{
	h(\tau+\de) = f(x_k-\tau \nabla f(x_k) - \de \nabla f(x_k)) = f(x_k-\tau\nabla f(x_k))-\de\dotp{\nabla f(x_k-\tau \nabla f(x_k))}{ \nabla f(x_k)} + o(\de).
}
Note that at $\tau=\tau_k$, $\nabla f(x_k-\tau \nabla f(x_k))=\nabla f(x_{k+1})$ by definition of $x_{k+1}$ in~\eqref{eq-grad-desc}.
 
If the current gradient is nonzero and the minimum is attained, an optimal step satisfies $\tau_k>0$, since small positive steps decrease $f$. The first-order condition is therefore
\eq{
	h'(\tau_k) = - \dotp{\nabla f(x_k)}{\nabla f(x_{k+1})} = 0.
}
Thus successive gradients, and hence successive descent directions, are orthogonal; see Figure~\ref{fig-gradesc}.
\end{rem}

\begin{rem}[Armijo rule]
	Instead of finding the optimal $\tau$, Armijo backtracking seeks an admissible $\tau$ that produces sufficient decrease.
	 
	Given $0<\al<1$ (values below $1/2$ are customary, especially for Newton-type methods), one considers a $\tau$ to be valid for a descent direction $d_k$ (for instance $d_k=-\nabla f(x_k)$) if it satisfies
	\eql{\label{eq-armijo-rule}
		f(x_k+\tau d_k) \leq f(x_k) + \al \tau \dotp{d_k}{\nabla f(x_k)}
	} 
	For small $\tau$, one has $f(x_k+\tau d_k) = f(x_k)+\tau\dotp{d_k}{\nabla f(x_k)}+o(\tau)$, so that, assuming $d_k$ is a valid descent direction (i.e $\dotp{d_k}{\nabla f(x_k)}<0$),  condition~\eqref{eq-armijo-rule} will always be satisfied for $\tau$ small enough (if $f$ is convex, the set of allowable $\tau$ is of the form $[0,\tau_{\max}]$).
	 
	Start from a positive trial step and repeatedly reduce it by $\tau\leftarrow\be\tau$, with $0<\be<1$, until~\eqref{eq-armijo-rule} holds.
	 
	This procedure is called backtracking line search.
\end{